% Created 2026-01-13 Tue 22:24
% Intended LaTeX compiler: lualatex
\documentclass[11pt]{article}
\usepackage{minted}
\usepackage{amsmath}
\usepackage{fontspec}
\usepackage{graphicx}
\usepackage{longtable}
\usepackage{wrapfig}
\usepackage{rotating}
\usepackage[normalem]{ulem}
\usepackage{capt-of}
\usepackage{hyperref}
\usepackage{minted}
\usepackage{amssymb}
\usepackage{xcolor}
\usepackage{newunicodechar}
\usepackage{newunicodechar}
\newunicodechar{₁}{\ensuremath{_1}}
\newunicodechar{₂}{\ensuremath{_2}}
\newunicodechar{₃}{\ensuremath{_3}}
\newunicodechar{₄}{\ensuremath{_4}}
\newunicodechar{₅}{\ensuremath{_5}}
\newunicodechar{₆}{\ensuremath{_6}}
\newunicodechar{₇}{\ensuremath{_7}}
\newunicodechar{₈}{\ensuremath{_8}}
\newunicodechar{₉}{\ensuremath{_9}}
\newunicodechar{₀}{\ensuremath{_0}}
\newunicodechar{ₐ}{\ensuremath{_a}}
\newunicodechar{ₑ}{\ensuremath{_e}}
\newunicodechar{ₒ}{\ensuremath{_o}}
\newunicodechar{ₓ}{\ensuremath{_x}}
\newunicodechar{ₙ}{\ensuremath{_n}}
\newunicodechar{ₖ}{\ensuremath{_k}}
\newunicodechar{ₘ}{\ensuremath{_m}}
\newunicodechar{≡}{\ensuremath{\equiv}}
\newunicodechar{…}{\ensuremath{\ldots}}
\newunicodechar{≈}{\ensuremath{\mathrel{\approx}}}
\newunicodechar{↓}{\ensuremath{\mathrel{\downarrow}}}
\newunicodechar{✅}{\textcolor{green!50!black}{\ensuremath{\checkmark}}}
\newunicodechar{❌}{\textcolor{red!70!black}{\ensuremath{\times}}}
\newunicodechar{∣}{\ensuremath{\mid}}
\newunicodechar{∈}{\ensuremath{\in}}
\newunicodechar{→}{\ensuremath{\to}}
\newunicodechar{⇒}{\ensuremath{\Rightarrow}}
\newunicodechar{↦}{\ensuremath{\mapsto}}
\newunicodechar{≥}{\ensuremath{\ge}}
\newunicodechar{≤}{\ensuremath{\le}}
\newunicodechar{≡}{\ensuremath{\equiv}}
\newunicodechar{⊂}{\ensuremath{\subset}}
\newunicodechar{⊃}{\ensuremath{\supset}}
\newunicodechar{⊆}{\ensuremath{\subseteq}}
\newunicodechar{⊇}{\ensuremath{\supseteq}}
\newunicodechar{⟶}{\ensuremath{\longrightarrow}}
\newunicodechar{⟵}{\ensuremath{\longleftarrow}}
\newunicodechar{⋯}{\ensuremath{\cdots}}
\newunicodechar{…}{\ensuremath{\ldots}}
\newunicodechar{∪}{\ensuremath{\cup}}
\newunicodechar{∩}{\ensuremath{\cap}}
\newunicodechar{∧}{\ensuremath{\wedge}}
\newunicodechar{∨}{\ensuremath{\vee}}
\newunicodechar{¬}{\ensuremath{\neg}}
\newunicodechar{∀}{\ensuremath{\forall}}
\newunicodechar{∃}{\ensuremath{\exists}}
\newunicodechar{⊢}{\ensuremath{\vdash}}
\newunicodechar{⊨}{\ensuremath{\vDash}}
\newunicodechar{⊥}{\ensuremath{\bot}}
\newunicodechar{⊤}{\ensuremath{\top}}
\newunicodechar{∅}{\ensuremath{\emptyset}}
\newunicodechar{↔}{\ensuremath{\leftrightarrow}}
\newunicodechar{⇔}{\ensuremath{\Leftrightarrow}}
\usepackage[
paper=letterpaper,
left=1in,
right=1in,
top=1in,
bottom=1in
]{geometry}
\usepackage{fontspec}
\setmainfont{Avenir Next}[
UprightFont = *,
ItalicFont = * Italic,
BoldFont = * Demi Bold,
BoldItalicFont = * Demi Bold Italic
]
\author{Michael Hohn}
\date{\today}
\title{}
\hypersetup{
 pdfauthor={Michael Hohn},
 pdftitle={},
 pdfkeywords={},
 pdfsubject={},
 pdfcreator={Emacs 30.2 (Org mode 9.7.11)}, 
 pdflang={English}}
\begin{document}

\tableofcontents

\section{MRVA deployment documentation}
\label{sec:orgfc0bd4b}
To understand what is being deployed we start with a brief overview of the MRVA
components and what they are.  The structure is illustrated at the container
level; the contents of the containers are described in full detail in their
\texttt{Dockerfiles}.  The full assembly sequence, from source, is described in the
\hyperref[sec:org6f1c38a]{Build Procedure}.  This deployment document stays at the container level, only
referencing the developer details.

The full list of containers is found in \url{../submodules/mrva-docker/docker-compose-demo.yml}
and their Dockerfiles in \url{../submodules/mrva-docker/containers/}.

The ones part of the mrva system are submodules of this repository --
\url{../submodules/}. They are
\begin{itemize}
\item gh-mrva
\item mrvaagent
\item mrvaserver
\item mrvacommander
\item mrvahepc
\end{itemize}

Also used are public containers:
\begin{itemize}
\item postgres
\item rabbitmq
\item minio
\end{itemize}

To clarify their roles and responsibilities, we can group these into
\begin{enumerate}
\item Internal containers:  MRVA proper
\item Stock containers: public containers providing services in the mrva network
\item Example service: This container is provided only as example service to make
MRVA functional.  The expectation is that this service is externally
provided in deployments
\end{enumerate}

With this grouping, we can describe the containers' function:
\begin{enumerate}
\item Internal containers:  MRVA proper, consisting of 
\begin{itemize}
\item \texttt{gh-mrva}.  A user/shell-driven service with three functions:
\begin{enumerate}
\item submit request
\item check status
\item download results (and optionally, DBs)
\end{enumerate}
\item \texttt{mrvaserver}.  The outside-facing service, it receives \texttt{gh-mrva} requests
and orchestrates their processesing by the other services.
\item \texttt{mrvaagent}.  This is the internal job runner; it performs performs the
analyses and returns results.  Specifically:
\begin{itemize}
\item pick up jobs
\item picks up DBs via / from from mrvahepc
\item run \texttt{codeql database analyze}
\item return results to rabbitmq
\end{itemize}
\item \texttt{mrvacommander}.  This is pure library code for \texttt{mrvaagent} and \texttt{mrvaserver}
\item mrva-docker -- docker containers assembled from other repos
\url{https://github.com/hohn/mrva-docker}
\end{itemize}

\item stock containers / services on deployment
\begin{itemize}
\item rabbitmq
\item postgres
\item minio
\end{itemize}

\item provided service / external.  This container is provided only as example;
the expectation is that this service is externally provided.

\begin{itemize}
\item mrvahepc -- endpoint for codeql DBs and metadata
\begin{itemize}
\item in-container-network service provides metadata
\item external service (here, https), provides the actual databases
\item this service intended for replacement

\item reference:
\begin{itemize}
\item \url{https://github.com/hohn/mrvahepc/tree/master}
\item \url{https://github.com/hohn/mrvahepc/tree/master?tab=readme-ov-file\#usage-sample-for-containers}
\item schema / json format: \url{https://github.com/hohn/mrvahepc/blob/master/bin/mc-hepc-init\#L65}
\end{itemize}

\item \textbf{Fixed API}:
\begin{itemize}
\item schema / json format: \url{https://github.com/hohn/mrvahepc/blob/master/bin/mc-hepc-init\#L65}
\item serves metadata -- provide info
\{"git\textsubscript{branch}": "HEAD", "git\textsubscript{commit}\textsubscript{id}": "dc0c7f3de40530053189c572936ae4fd1567269b", "git\textsubscript{repo}": "bulk-builder", "ingestion\textsubscript{datetime}\textsubscript{utc}": "2022-11-30T00:44:34.455172576Z", "result\textsubscript{url}": "\url{http://hepc/db/db-collection.tmp/bulk-builder-bulk-builder-ctsj-a8e72a.zip}", "tool\textsubscript{id}": "9f2f9642-febb-4435-9204-fb50bbd43de4", "tool\textsubscript{name}": "codeql-cpp", "tool\textsubscript{version}": "2.11.3", "projname": "bulk-builder/bulk-builder"\}
\item served content listed in metadata -- via https in this example.
The field is fixed, the url is arbitrary.
\begin{itemize}
\item "result\textsubscript{url}":
"\url{http://hepc/db/db-collection.tmp/bulk-builder-bulk-builder-ctsj-a8e72a.zip}"
\end{itemize}
\end{itemize}
\end{itemize}
\end{itemize}
\end{enumerate}

The container assembly is driven via the \texttt{mrva-docker} repository.
\section{Build Procedure}
\label{sec:org6f1c38a}
\end{document}
